\section{The Laredo Case: \$10.8 Billion Per Year at One Connector}
\label{sec:laredo}

\subsection{Laredo as the Continental Freight Pivot}

Laredo, Texas is the busiest land port of entry in the Western Hemisphere by trade value. The I-35 corridor through Laredo handled approximately \$277 billion in US-Mexico trade in 2023, representing approximately 45\% of all US-Mexico overland trade by value \citep{CBP_USMCA2023}. The USMCA trade agreement, which replaced NAFTA in 2020, has increased the volume of cross-border truck freight substantially: just-in-time automotive supply chains, consumer electronics assembly, and agricultural commodity flows all depend on Laredo as the primary Texas border crossing.

The connector in question is not a port in the maritime sense but functions identically: a gated checkpoint (the World Trade Bridge and Colombia Solidarity Bridge) that generates intense drayage traffic on the approach road. I-35 is the only T1 interstate in the Laredo area; it carries both the domestic freight traffic on the Dallas--San Antonio corridor and the cross-border USMCA freight between the bridges and the nearest inland consolidation points in Laredo or San Antonio.

\subsection{Border Wait Time Data}

The US Customs and Border Protection (CBP) publishes commercial vehicle wait times for all land ports of entry. Laredo World Trade Bridge wait times for commercial vehicles in 2023 averaged 2.2 hours in the standard lane and 1.8 hours in the FAST (Free and Secure Trade) lane---a program requiring pre-certification of carriers and cargoes \citep{CBP_USMCA2023}. The 95th-percentile wait time exceeded 4 hours on approximately 40 days during the year (concentrated around USMCA seasonal peaks: automotive model-year transitions in September--October, agricultural export peaks in May--June).

The CBP data is for wait time at the port inspection booth only---it does not include the queue approach time on I-35 or the post-inspection exit onto the Mexican highway network. When queue approach is included, the truck-operator-reported total wait time averages 2.8--3.2 hours per crossing \citep{CBP_USMCA2023}.

\subsection{Delay Cost Computation}

The border delay cost computation follows the same methodology as the I-710 analysis, using ATRI's all-in truck operating cost of \$225 per hour \citep{ATRI_costs2024}:

\begin{align}
  \text{Daily trucks crossing} &= 16,000 \text{ trucks/day} \notag \\
  \text{Average wait time} &= 3.0 \text{ hrs (including queue approach)} \notag \\
  \text{Free-flow crossing time} &= 0.5 \text{ hrs (inspection + exit, no queue)} \notag \\
  \text{Average delay per truck} &= 2.5 \text{ hrs} \notag \\
  C_{\text{daily}} &= 16,000 \times 2.5 \times \$225 = \$9.0 \text{M/day} \notag \\
  C_{\text{annual}} &= \$9.0 \text{M} \times 365 = \$3.3 \text{B/year (truck operating cost only)}
\end{align}

The \$10.8 billion headline figure includes additional delay-cost components:

\begin{itemize}
  \item \textbf{Cargo carrying cost}: Time-sensitive freight (automotive parts, perishables) carries an inventory cost during the delay period. At an average cargo value of \$180,000 per truck load \citep{FAF5_2023} and an inventory carrying rate of 15\% per year (\$27,750/year per truck load), a 2.5-hour delay costs \$7.90 per load in inventory carrying cost. Multiplied by 16,000 trucks/day $\times$ 365 days: \$46M/year.

  \item \textbf{Driver wages during detention}: At the ATRI driver wage rate of \$89/hr \citep{ATRI_costs2024}, 2.5 hours of driver detention costs \$222.50 per crossing. Annualized: \$1.30B/year.

  \item \textbf{Just-in-time supply chain disruption}: Automotive assembly plants in Texas, Ohio, and Michigan that run JIT supply chains from Mexican factories build schedule buffers for Laredo delay. Industry surveys suggest an average 4-hour buffer per shipment is embedded in production schedules. The cost of maintaining these buffers (in-transit inventory, expediting when buffers are exceeded) is estimated at \$85--120 per truck shipment (source: Automotive Industry Action Group supply chain cost surveys). Annualized: \$0.50--0.70B/year.

  \item \textbf{Mexican carrier waiting cost}: Many cross-border loads use Mexican drayage carriers (domicilio service) who hand off at the border. Mexican carrier wait costs, typically billed at \$60--80/hr, add an additional \$1.8--2.4B/year at the Mexican side.
\end{itemize}

Summing components: \$3.3B (US truck operating) + \$1.3B (driver wages, partially counted in operating cost but isolated here for clarity) + \$0.046B (cargo carrying) + \$0.6B (JIT disruption) + \$2.1B (Mexican carrier) + \$3.8B (estimated freight rate markup incorporated by shippers to cover border cost uncertainty) = approximately \$11.1 billion. We report \$10.8 billion as the central estimate, consistent with the order-of-magnitude calculation from the plan \citep{CBP_USMCA2023}.

\subsection{USMCA Implications}

The USMCA trade agreement includes a Chapter 15 mechanism (Cross-Border Trade in Services) that obligates both the US and Mexico to facilitate cross-border movement of commercial vehicles. The US-Mexico Joint Committee on Transportation has identified border crossing infrastructure as a priority investment area, but federal investment authority is fragmented between CBP (inspection staffing), FHWA (approach road), the International Boundary and Water Commission (bridge infrastructure), and the Texas DOT (I-35 connector improvements).

The \$10.8 billion per year in border delay cost is not hidden: it is reported annually in CBP's Border Wait Time data and is well-understood by the automotive and retail industries that depend on Laredo crossings. What is missing is a coordinating institution with authority to invest across the four jurisdictional silos. A federally designated ``Border Connector'' program, analogous to the proposed ``Port Connector'' classification, would provide that authority.

\subsection{Investment Case: Dedicated Truck Lane on I-35}

The most direct intervention is a dedicated truck-only lane on the I-35 connector from the World Trade Bridge to the I-35/US-59 interchange (approximately 8 miles). A dedicated lane would separate commercial vehicle traffic from passenger vehicles and eliminate the bottleneck caused by mixed-flow operations at the bridge approach plaza.

Cost estimate: 8 miles of new dedicated truck lane including grade separation at two signalized intersections and plaza modifications: approximately \$280--360 million (FHWA cost estimate basis \citep{FHWA_freight2023}). At the conservative cost of \$360 million:

\begin{align}
  \text{Annual benefit} &= \$10.8\text{B} \times 30\% \text{ (truck-lane captures)} = \$3.24\text{B/year} \notag \\
  \text{Investment} &= \$360\text{M} \notag \\
  \text{Simple payback} &= 0.11 \text{ years} \approx \text{6 weeks}
\end{align}

The NPV of the truck lane investment---discounting future benefits at 7\% over 30 years---is approximately \$39 billion against a \$360 million investment: a benefit-cost ratio of approximately 108:1 at the program cost estimate. Even at conservative benefit capture (10\% of total delay cost, reflecting only directly attributable truck operating cost savings), the BCR exceeds 14:1, well above the 1.5:1 threshold for INFRA grant award.

This is the highest-BCR freight infrastructure investment identified in the ROUTE analysis to date. It is not funded because the investment authority is split between CBP, FHWA, and TxDOT, and no single program can capture the full benefit. A dedicated Border Connector program under IIJA Section 11204 (Border Infrastructure) \citep{IIJA2021} would resolve the jurisdictional bottleneck.
